\section{Definitions}
\label{sec:definitions}

We state the definitions explicitly so that the falsifiability conditions are exposed.

\paragraph{Setting.}
Let $\mathcal{M}$ be a transformer (Gemma-2-2B-it), $\mathcal{D}$ a CLT-HP feature dictionary~\cite{hanna2025clthp,ameisen2025circuittracing}, and $G_p = (V_p, E_p)$ the attribution graph for prompt $p$ with target logit $y$. Each node $f \in V_p$ is a (feature, position) pair with influence $\mathrm{infl}(f) \in \mathbb{R}_{\geq 0}$; the cumulative-influence threshold $\tau$ selects a feature universe $V_p^{\tau} \subseteq V_p$.

\begin{definition}[Supernode]
\label{def:supernode}
A \emph{supernode} $s = (\rho, \nu, F)$ is a triple of a functional role $\rho \in \{\textsc{Sem-Dict}, \textsc{Sem-Conc}, \textsc{Rel}, \textsc{Say-X}\}$, a string label $\nu$, and a feature set $F \subseteq V_p^{\tau}$ all sharing role $\rho$ and label $\nu$. Each feature belongs to at most one supernode (\cref{sec:method}).
\end{definition}

\begin{definition}[Concept-aligned supernode]
\label{def:aligned}
A supernode $s = (\rho, \nu, F)$ is \emph{aligned to concept $c$} if (i) $\nu$ matches $c$ under the substring-aware matcher of \cref{sec:method}, or (ii) $\rho = \textsc{Say-X}$ and the discovered target token of $s$ matches $c$. Each entity $e$ in a domain induces a set of concepts $C(e)$ across its semantic fields (input, intermediate, answer); the corresponding supernodes are $S(e) = \{s : s \text{ aligned to some } c \in C(e)\}$.
\end{definition}

\begin{definition}[Additive entity-swap intervention]
\label{def:swap}
For a swap from source entity $e_A$ to target entity $e_B$ on prompt $p$, the intervention $\mathcal{I}(e_A \to e_B; M_a, M_b)$ ablates each feature $f \in \bigcup_{s \in S(e_A)} F$ by additive decoder injection with multiplier $M_a$ (default $-2$) and amplifies each $f \in \bigcup_{s \in S(e_B)} F$ with $M_b$ (default $20$). Attention is \emph{not} frozen during intervention ($\textsc{freeze\_attention}{=}\texttt{false}$), so the model can route around the perturbation~\cite{ameisen2025circuittracing}.
\end{definition}

\begin{definition}[Hit and vsMax]
\label{def:hit}
Let $g(p, \mathcal{I})$ be the generated continuation under intervention $\mathcal{I}$ at temperature $0.3$, $n_{\mathrm{tokens}}{=}10$, frequency penalty $2.0$, seed $42$. A swap is a \textbf{Hit} for target $e_B$ iff the first-subword token of $e_B$'s answer appears in $g$. The metric \textbf{vsMax} is $\max_{t \leq 10}\big(\ell_{e_B}(t) - \max_{e' \in \mathcal{A}\setminus\{e_B\}} \ell_{e'}(t)\big)$, where $\ell_e(t)$ is the logit of $e$'s first-subword token at trajectory position $t$ and $\mathcal{A}$ is the domain answer set.
\end{definition}

\begin{definition}[Matched random control]
\label{def:matched}
The \emph{matched random control} for swap pair $(e_A, e_B)$ is an intervention $\mathcal{I}^{\mathrm{rand}}$ identical to $\mathcal{I}(e_A \to e_B)$ in feature count and per-layer histogram, with features sampled uniformly at random from $V_p^{\tau} \setminus \bigcup_{e \in \mathcal{E}} \bigcup_{s \in S(e)} F$ (i.e., outside \emph{every} concept-aligned supernode in the domain). Three deterministic replicates are drawn per pair, seeded by $\mathrm{sha256}(\mathrm{run\_seed} : \mathrm{pair\_id} : \mathrm{rep} : \mathrm{mode})$.
\end{definition}

\begin{definition}[Operationally useful label]
\label{def:opuseful}
The supernode label $\nu$ for entity $e$ is \emph{operationally useful} on prompt family $p$ iff both conditions hold over all swap pairs $(e_A, e_B)$ in the domain:
\begin{itemize}\itemsep0pt
\item \textbf{Condition 1 (Specificity).} The Hit-rate of $\mathcal{I}(e_A \to e_B)$ exceeds $5$~pp on a per-pair basis, and labeled vsMax is positive on average.
\item \textbf{Condition 2 (Matched-Control).} Across replicates, $\mathcal{I}^{\mathrm{rand}}$ achieves Hit-rate at least $5$~pp lower than the labeled intervention and labeled vsMax exceeds matched-control vsMax by at least $1.0$ logit.
\end{itemize}
A label fails to be operationally useful if either condition fails; this is the falsification criterion used throughout the paper.
\end{definition}
